\lecture{21}{Wed 29 Nov 2023 10:40}{}

\begin{proposition}
	Let $M_t$ be a continuous $L^2$-martingale. Then it has bounded variation only when it is constant, i.e., $M_t \equiv M_0$.
\end{proposition}
\begin{proof}
If $M _t$ is bounded variation on $[0,t]$, then $\left\langle M \right\rangle _t = 0$. To verify, note that
\[
	\sum_{k = 0}^{n - 1} \left| M_{t_{k+1}} - M_{t_k} \right| ^2
	\le \sup _{0\le k\le n-1} \left| M_{t_{k+1}} - M_{t_k} \right| \sum_{k=0}^{n-1} \left| M_{t_{k+1}} - M_{t_k} \right|
.\] 
Where the $\sup $ goes to zero by continuity and the second term goes to the total variation which is bounded.

Therefore, $M_t^2$ is a martingale. We also know that $M_t$ is a martingale. Then
\[
	E\left( \left( M_t - M_s \right) ^2 | \mathcal{F}_s\right)	
	= E\left( M_t^2 - M_s^2 | \mathcal{F}_s \right) 
	= 0
.\] 
Take expectation again,
\[
	E | M_t - M_s| ^2 = 0 \implies M_t \text{ is constant}
.\] 
%This can be shown from Chebyshev: $P\left( M_t - M_s > \varepsilon \right) \le \frac{E\left( M_t - M_s \right) ^2}{\varepsilon^2} = 0$.
\end{proof}

Thus, we can proceed exactly as in the case of BM.

\begin{definition}[Ito integral](Elementary functions, general continuous martingale)
	
Consider $M$ a continuous $L^2$ martingale.
Let $L_M^2$ be space of progressively measurable processes $u_ t$ wrt $\mathcal{F}_t$ such that
\[
	E\left( \int _0^t u_s^2 d\left\langle M \right\rangle _s \right) < \infty 
.\] 
Define the elementary processes $\mathcal{E}$ by
\[
	\mathcal{E} = \left\{ u_t \in L^2_M: u_t = \sum_{k = 0}^{n - 1} F_k 1_{[t_k, t_{k+1})}(t), E|F_k|^2 < \infty , F_k \in \mathcal{F}_{t_k} \right\} 
.\] 

Let $u \in \mathcal{E}$ be elementary, taking value $F_k $ on $[t_k, t_{k+1})$, then define
\[
	I_M(u) = \int _0^T u_s d M_s = \sum_{k = 0}^{n - 1} F_k \left( M_{t_{k+1}} - M_{t_k} \right) 
.\] 
\end{definition}

Proceed in the same way as previous theorem about Brownian motion, we have

\begin{theorem}[Ito Integral]
	There exists a unique linear map $I_M: L^2_M \to  L^2_M$ such that
	\begin{enumerate}
		\item $I_M(u) = \sum_{k = 0}^{n - 1} F_k \left( M_{t_{k+1}} - M_{t_k} \right) , u \in \mathcal{E}$ (Ito integral for elementary functions)
		\item $E\left( I_M(u) \right) ^2 = E\left( \int _0^T u_s^2 d \left\langle M \right\rangle _s \right) , u \in \mathcal{E}$ (Ito isometry for elementary functions)
		\item $\mathcal{E}$ is dense in $L^2_M$
		\item  $I_M(u) =_{L^2} \lim_{m \to \infty} I_M(e^m), e^m \in \mathcal{E}$ (Continuity of $I_M$)
		\item $E\left( I_M(u) \right) ^2 = E\left( \int _0^T u_s^2 d \left\langle M \right\rangle _s \right) $, $u \in L^2_M$ (Ito isometry)
	\end{enumerate}
\end{theorem}
The proof is completely parallel to the $M_t = B_t$ case.

\begin{remark}
	We have $\int _0^t u_s d M_s$ is a continuous martingale (equivalent of Theorem~\ref{thm:ito-martingale-continuous}).
	Applying Ito isometry and consider the continuous martingale $\left( \int _0^t u_s d M_s \right) ^2 - \left\langle \int _0^t u_s d M_s \right\rangle $, we have
	\[
		\left\langle \int _0^\cdot u_s d M_s \right\rangle _t = \int _0^t u_s^2 d \left\langle M \right\rangle _s
	.\] 
\end{remark}

\begin{example}
	As for BM,  $\int _0^t u_s d B_s$ is a martingale, and in fact it is continuous.
\end{example}
\begin{example}
	When $M_t  = \int _0^t H_s d B_s$. Then
	\[
		\left\langle M \right\rangle _t = \int _0^t |H_s|^2 ds
	.\] 
	We have $d \left\langle M \right\rangle _t = |H_t|^2 dt$. Thus we have
	\[
		u \in L_M^2 \iff  
		E\left( \int _0^t |u_s|^2 d \left\langle M \right\rangle _s \right)  = E \left( \int _0^t |H_s|^2 u_s^2 ds \right)  < \infty 
	.\] 
	Thus $u_s H_s \in L_B^2$. 
	In this case we have,
	\[
		\int _0^t u_s d M_s = \int _0^t u_s H_s \cdot d B_s
	.\] 
	To justify, both are continuous martingales, and their quadratic variation are identical. Continuous martingales are characterized by their quadratic variations, as we will verify in the Martingale Representation Theorem.
\end{example}

\begin{remark}
	\begin{enumerate}
		\item In order to define stochastic integral, we only need local integrability on the martingale in time, i.e. $E M_t^2 < \infty $ for all $t$.

\begin{definition}
	A stochastic process $M_t$ adapted to $\mathcal{F}_t$ is called a \textit{local martingale} if there exists stopping times $T_n$ going to infinity, such that $M_{t \wedge T_n}$ is a martingale.
\end{definition}

\item A \textit{semimartingale} is a process of the form
	\[
		X_t = A_t + M_t
	.\] 
	where $M_t$ is a local martingale, and $A_t$ is a process of locally finite variation. We can define
	\[
		\int _0^t H_s d X_s = \int _0^t H_s d A_s + \int _0^t H_s d M_s
	.\] 
	where the first integral is just Riemann-Stieltjes integral, and the second is defined for local martingales.
	\end{enumerate}
\end{remark}

\begin{lemma}
	Let $X_t = X_0 + A_t + M_t$ be a continuous semimartingale. Let $\left| \Delta_n(t) \right|  \to  0$ as $n \to \infty $, then
	\[
		\lim \sum_{n=0}^{n-1} \left| X_{t_{n+1}} - X_{t_n} \right| ^2 = \left\langle M \right\rangle _t
	.\] 
	i.e., $\left\langle X \right\rangle _t = \left\langle M \right\rangle _t$
\end{lemma}
\begin{proof}
	\begin{align*}
		\left| X_{t_{k+1}} - X_{t_k} \right| ^2
		&= \left| (A_{t_{k+1}} - A_{t_k})+(M_{t_{k+1}} - M_{t_k}) \right|^2  \\
		&= \left| A_{t_{k+1}} - A_{t_k}\right|^2
		+ 2(A_{t_{k+1}} - A_{t_k}) (M_{t_{k+1}} - M_{t_k})
		+ \left| M_{t_{k+1}} - M_{t_k}\right|^2  
	.\end{align*}
	The first two terms go to zero when we take sum and let $\Delta_k \to 0 $, using the fact that $A$ is bounded variation on $[0,t]$ and $M$ is continuous.
\end{proof}
\begin{remark}
	We have same argument for $\left\langle X, Y \right\rangle _t = \left\langle M_t, N_t \right\rangle $ where $X_t = A_t + M_t$, $Y_t = C_t + N_t$.
\end{remark}

Furthermore,
\begin{lemma}
	For $f, g \in C^2$, $X, Y$ semimartingales,
	\[
		\left\langle f(X), g(Y) \right\rangle 
		= \int _0^t f'(X_s) g'(Y_s) d \left\langle X, Y \right\rangle _s
	.\]
\end{lemma}

\subsection{Ito Formula}

\subsubsection{Motivation and special case}

Motivation comes from integration by parts. If we know $I(X)$ and $I(Y)$, we can know $I(f(X,Y))$ where $f$ is polynomial. Now we want to generalize to all $C^2$ functions.
